% !TeX encoding = UTF-8 Unicode

\DocumentMetadata{lang=zh-CN,pdfversion=1.7}
\documentclass[a4paper,twoside,onecolumn,12pt,fontset=none,UTF8,notitlepage]{ctexrep}
\usepackage{fancyhdr,titlesec}%标题、脚注格式
\usepackage{luacolor, xcolor}%颜色
\usepackage[soul]{lua-ul}%下划线与其他装饰
\usepackage{graphicx}%插入图片
\usepackage{float}%图表位置，如表格过长致横向排版已不可能，应加用adjustboxs宏包排版旋转表格
\usepackage{tabularray,multicol}%表格排版
\usepackage{enumitem}%列表排版
\usepackage[a4paper, hmargin=1.5cm, top=1.5cm, bottom=2cm, headheight=0pt, ignorehead]{geometry}%页边距与排版
\usepackage[bookmarksnumbered=true, hidelinks]{hyperref}%超链接
\usepackage{texosquery}%获取秒数及时间
\usepackage{datetime2}%扉页基本时间及格式
%\usepackage{showframe}%显示页面布局与边框，仅Preview版本排查错误

% 设置文档标题
\hypersetup{%
    pdftitle = {山二医指南（部分） SDSMU-Companion-Part},%
    pdfauthor = {LinkChou, etc.},%
    pdfdisplaydoctitle = true%
}

% 重设定标题与正文间距
\ctexset{
    chapter={
        beforeskip=0ex,
        afterskip=4ex
    },
    section={
        beforeskip=0ex,
        afterskip=0ex
    },
    subsection={
        beforeskip=0ex,
        afterskip=0ex
    },
    subsubsection={
        beforeskip=0ex,
        afterskip=0ex
    }
}

% 重定义文章字体为梦源宋体（不使用思源宋体的原因是思源的行高过高）
\setmainfont[BoldFont=Dream Han Serif CN W20.ttf]{Dream Han Serif CN W7.ttf}
\setCJKmainfont{Dream Han Serif CN W7.ttf}[BoldFont=Dream Han Serif CN W20.ttf]
\setCJKsansfont[BoldFont=Dream Han Serif CN W20.ttf]{Dream Han Serif CN W7.ttf}
\setCJKmonofont[BoldFont=Dream Han Serif CN W20.ttf]{Dream Han Serif CN W7.ttf}

% 段间距与行间距重设定
\setlength{\parskip}{.75ex plus 5pt minus 8pt}
\setlength{\lineskip}{.5ex plus 5pt minus 8pt}

% 页眉页脚
\pagestyle{fancy}
\fancyhf{}
\fancyfoot[C]{\thepage}

% 脚注格式（无分割线）
\renewcommand{\headrulewidth}{0pt}
\renewcommand{\footrulewidth}{0pt}
\renewcommand{\thefootnote}{\fnsymbol{footnote}}

% 列表样式重定义
\setlist[enumerate,1]{
    label=\arabic*.,
    ref=\thesubsubsection-\theenumi
}
\setlist[enumerate,2]{
    label=(\arabic*),
    ref=\thesubsubsection-\theenumi.\theenumii
}
\setlist[enumerate,3]{
    label=[\arabic*],
    ref=\thesubsubsection-\theenumi.\theenumii.\theenumiii
}
\setlist[enumerate,4]{
    label=<\arabic*>,
    ref=\thesubsubsection-\theenumi.\theenumii.\theenumiii.\theenumiv
}
\setlist[enumerate]{
    itemsep=0pt
}

%重定义表格脚注字体大小并增加“no-caption”表格样式
\DefTblrTemplate{note-text}{default}{{\footnotesize\InsertTblrNoteText}}
\DefTblrTemplate{remark-tag}{default}{{\footnotesize\bfseries\InsertTblrRemarkTag}}
\DefTblrTemplate{remark-text}{default}{{\footnotesize\InsertTblrRemarkText}}
\NewTblrTheme{no-caption}{
    \SetTblrTemplate{caption}{empty}
}

%重定义表格表注的表现形式避免问题
\renewcommand\TblrOverlap[1]{#1}

%重定义表格换页时内容
\DefTblrTemplate{contfoot-text}{default}{续下页}
\DefTblrTemplate{conthead-text}{default}{（续）}

% 重定义超链接与引用表现样式
\newcommand{\uhref}[2]{%
    \ul{\href{#1}{#2}}%
}
\newcommand{\uref}[1]{%
    \ul{\ref{#1}}%
}

% 定义下划线样式
\newcommand{\uul}[1]{%
  \ul[bottom=\dimexpr-.7ex+1pt\relax, height=.1ex]{%
    \ul[top=-.7ex, height=.1ex]{%
      #1%
    }%
  }%
}

% 设置校徽命令
\newcommand{\BackgroundPic}{
    \noindent\parbox{\paperwidth}{
        \centering
        \vspace*{200pt}
        %\vspace*{350pt}
        \includegraphics[width=5cm]{logo.pdf}
    }
}

% 颜色重设定
\newcommand{\colored}[4]{
    \textcolor[RGB]{#1,#2,#3}{#4}
}

% 时间样式重设定
\newcommand{\CurrentCustomTime}{
    \DTMsavenow{Compiled_time}%保存编译时间
    \DTMfetchyear{Compiled_time}.%
    \DTMtwodigits{\DTMfetchmonth{Compiled_time}}.%
    \DTMtwodigits{\DTMfetchday{Compiled_time}}\quad%
    \DTMtwodigits{\DTMfetchhour{Compiled_time}}:%
    \DTMtwodigits{\DTMfetchminute{Compiled_time}}:%
    \DTMtwodigits{\DTMfetchsecond{Compiled_time}}\quad%
    GMT\DTMfetchTZhour{Compiled_time}:\DTMfetchTZminute{Compiled_time}
}

\begin{document}

% 扉页

% 添加标题校徽背景图
\AddToHookNext{shipout/background}{\BackgroundPic}

% 题目
\title{%
    \vspace{100pt}
    {\Huge\textbf{山东第二医科大学\\[5pt]——？？？}}
}
\author{%
    版本：\textbf{1.0} \qquad%
    CC BY-SA 4.0 LICENSE \qquad%
    \uhref{Mailto:LinkChou@yandex.com}{LinkChou}}
\date{\CurrentCustomTime}
\maketitle

% 脚注样式重设定
\renewcommand{\thefootnote}{\arabic{footnote}}

\vspace{100pt}

\part[敬告]{敬告}
力薄才疏，不免舛误，敬请绳愆纠缪。具体内容均以学校官方为准，如有变更恕不另行通知。

\part[总述]{总述}
公交运营时间通常为06:30～19:00

免费线路\footnotemark：9、13、19、29、63、69、71、101、109、166、167
\footnotetext{据《滨海及浮烟山两地大学城在校学生免费乘坐校区至中心城区公交车实施方案》（潍交城〔2024〕5号），政策暂行时间为：2024.4.29～2025.4.28。2025.08.16更新，因新的政策尚未下发，现在本指南仅可作为公交线路参考。}

在\textbf{填表报名}并向学校\textbf{提交相关信息}后，可通过\textbf{与提交信息相符的手机号}注册认证『潍坊公交大学生乘车码』微信小程序免费乘坐上述车次，新生可持学生证或录取通知书临时代替。如需更改绑定的手机号，需向学校提出申请或在每年下发的《免费乘车大学生基础信息表》中按要求注明。

\part[常用站点名称表]{常用站点名称表}
\begin{tblr}[
        long,
        caption = {常用站点名称对应关系一览表},
    ]{
        cells = {c,m},
        rowhead = 1,
        row{1} = {cmd=\bfseries},
        vlines,
        hlines,
        cell{2,8,30,32}{1} = {r=2}{},
        cell{10,20,23}{1} = {r=3}{},
        cell{4,13}{1} = {r=4}{},
        cell{28,31}{3} = {r=2}{},
        hline{1-2,8,13,17,Z} = {-}{1pt},
    }
    %尽量手动控制，尽量避免因自动断页导致的underfull vbox警告
    常见目的地       & 具体方位     & 站点名                   \\
    浮烟山校区       & 南门         & 山东第二医科大学南门     \\*
                     & 北门         & 山东第二医科大学北门     \\\nopagebreak[2]
    虞河校区         & 北门         & 胜利街虞河路路口西       \\*
                     & 东北侧       & 虞河路胜利街路口北       \\*
                     & 东门         & 虞河路胜利街路口南       \\*
                     & 西侧         & 鸢飞路胜利街路口北       \\
    附属医院         & 西门         & 附属医院                 \\*
                     & 南侧         & 福寿街虞河路路口东       \\
    人民医院         & 本部东门     & 人民医院东门             \\*
                     & 本部西侧     & 人民医院                 \\*
                     & 北辰院区南门 & 北辰院区                 \\
    火车站           & 北广场       & 火车站                   \\*
                     & 东侧         & 青年路铁路桥北           \\*
                     & 东北侧       & 火车站                   \\*
                     & 潍坊北站     & 高铁北站                 \\
    人民公园         & 西门         & 人民公园西门             \\
    市中医院         & 东门         & 市中医院                 \\
    妇幼保健院       & 东门         & 潍坊市妇幼保健院潍城院区 \\
    泰华城           & 北侧         & 泰华城                   \\*
                     & 西南侧       & 青年路胜利街路口南       \\*
                     & 南侧         & 风筝广场                 \\
    万达广场         & 西侧         & 鸢飞路福寿街路口南       \\*
                     & 南侧         & 潍坊中学                 \\*
                     & 北侧         & 福寿街鸢飞路路口东       \\
    谷德广场         & 南侧         & 谷德广场                 \\
    奎文防疫站       & 无           & 奎文防疫站               \\
    鲁台会展中心     & 北门         & 鲁台会展中心             \\*
    蓝海大饭店       & 南门         &                          \\*
    谷德锦           & 南侧         & 北宫街清平路路口西       \\*
                     & 东侧         & 潍坊交通职业中等专业学校 \\*
    奥林匹克体育公园 & 东侧         &                          \\*
                     & 南侧         & 奥体中心                 \\
    十笏园           & 西侧         & 十笏园
\end{tblr}

\part[常用目的地表]{常用目的地表}
\begin{tblr}[
        long,
        caption = {常用路线汇总表},
        note{1} = {加粗线路为上述免费乘车线路。},
    ]{
        rowhead = 1,
        cells = {c,m},
        row{1} = {cmd=\bfseries},
        cell{1}{2} = {c=3}{},
        cell{2,7,9,11,15,20,26}{1} = {r=2}{},
        cell{4}{1} = {r=3}{},
        cell{4,18,19,22,24,28}{1} = {}{cmd=\bfseries},
        vlines,
        hlines,
        hline{1-2,Z} = {-}{1pt},
    }
    路线\TblrNote{1} & 常见目的地         &                          &                  \\*
    5                & 火车站（北广场）   & 市中医院（东门）         & 风筝广场         \\*
                     & 万达广场（西侧）   &                          &                  \\
    13               & 浮烟山校区（南门） & 虞河校区（北门/东北侧）  & 火车站（东侧）   \\*
                     & 附院（西门）       & 泰华城（南侧）           & 妇幼保健院       \\*
                     & 奎文防疫站         & 人民公园                 & 风筝广场         \\
    22               & 泰华城（北侧）     & 十笏园                   & 谷德锦           \\*
                     & 奥林匹克公园       &                          &                  \\
    23               & 火车站（北广场）   & 附院（南侧）             & 十笏园           \\*
                     & 万达广场（北侧）   &                          &                  \\
    25               & 虞河校区（东北侧） & 附院（西门）             & 火车站（东北侧） \\*
                     & 火车站（潍坊北站） & 奎文防疫站               & 妇幼保健院       \\
    D28              & 人民医院（东门）   & 人民医院（北辰院区南门） &                  \\
    56               & 火车站（北广场）   & 泰华城（北侧）           & 谷德广场         \\
    59               & 虞河校区（北门）   & 虞河校区（东门）         & 人民医院（东门） \\*
                     & 谷德锦（南侧）     & 奥林匹克公园（东南侧）   &                  \\
    66               & 虞河校区（北门）   & 市中医院                 & 万达广场（南侧） \\
    69               & 浮烟山校区（北门） & 火车站（东侧）           & 火车站（北广场） \\*
    71               & 浮烟山校区（北门） & 火车站（东侧）           & 火车站（北广场） \\
    75               & 虞河校区（西侧）   & 万达广场（西侧）         & 风筝广场         \\*
                     & 火车站（北广场）   &                          &                  \\
    101              & 浮烟山校区（北门） & 火车站（北广场）         &                  \\
    106              & 火车站（东北侧）   & 火车站（潍坊北站）       &                  \\
    109              & 浮烟山校区（南门） & 泰华城（东北侧）         &                  \\
    132              & 人民医院（西侧）   & 万达广场（西侧）         &                  \\
    162              & 虞河校区（东门）   & 人民医院（东门）         & 附院（西门）     \\*
                     & 奎文防疫站         &                          &                  \\
    167              & 浮烟山校区（南门） & 鲁台会展中心（北门）     & 蓝海大饭店
\end{tblr}

\end{document}
